\documentclass[11pt,a4paper]{article}
\usepackage[margin=2.5cm]{geometry}
\pagestyle{empty}
\hyphenpenalty=10000 \exhyphenpenalty=10000
\input{ocrbench}
\begin{document}
\ocrpar{Percolation on a lattice keeps each edge open with probability $p$ and asks whether the open cluster of the origin is infinite. Kesten proved that the critical value for the square lattice is one half, confirming an argument of Harris and a prediction from physics.}
\ocrpar{The theorem of K\"onig on infinite trees, the subadditive lemma of Fekete, and Kolmogorov's zero one law are the classical tools. G\"odel's name appears only in a historical remark, and the inequality of Harris and Kleitman, often called the FKG inequality, does most of the work.}
\ocrpar{Near the critical point the probability that the origin is connected to distance $n$ decays like a power of $n$. On the triangular lattice the exponent is known exactly, thanks to the conformal invariance proved by Smirnov and the stochastic Loewner evolution introduced by Schramm.}
\ocrpar{We prove a sharp threshold for the crossing of a long rectangle by an open path. The proof combines the Russo formula with an influence bound in the style of Kahn, Kalai and Linial, and it applies to every graph with enough symmetry. M\"obius inversion and the Ces\`aro mean appear in the appendix, together with a remark on the work of Pr\'ekopa.}
\end{document}
